\section*{Chapitre \ref{ch:b3:supramolecular} --- Chimie supramoléculaire}

\begin{solution}{exo:b3:supramolecular:1}
Ion--dipôle~; empilement $\pi$~; cation--$\pi$~; \omterm{def:b3:supramolecular:interactions}{liaison halogène} (de I vers N)~; trois liaisons hydrogène.
\end{solution}

\begin{solution}{exo:b3:supramolecular:2}
$\Delta G^\circ = -RT\ln K = -8.314 \times 298 \times \ln(\num{1.0e4}) = \qty{-22.8}{kJ/mol}$.
\end{solution}

\begin{solution}{exo:b3:supramolecular:3}
$K_{\ce{K+}}/K_{\ce{Na+}} = 10^{1.7} = 50$~; la différence vaut $RT\ln 50 = \qty{9.7}{kJ/mol}$ en faveur du potassium.
\end{solution}

\begin{solution}{exo:b3:supramolecular:4}
$n/(n + m) = 1/3$~: un \omterm{def:b3:supramolecular:supramolecular}{hôte} pour deux \omterm{def:b3:supramolecular:supramolecular}{invités}, $\mathrm{HG_2}$.
\end{solution}

\begin{solution}{exo:b3:supramolecular:5}
$H_0 + G_0 + 1/K = \qty{4.667e-3}{mol/L}$~; $[\mathrm{HG}] = (4.667 - \sqrt{4.667^2 - 4 \times 2.0 \times 2.0})/2 \times 10^{-3} = \qty{1.13e-3}{mol/L}$~: fraction
liée 0.566, $\delta = 3.560 + 0.160 \times 0.566 = \qty{3.651}{ppm}$.
\end{solution}

\begin{solution}{exo:b3:supramolecular:6}
$\log K = 18.3 - 8.6 = 9.7$, $K = \num{5e9}$~; $\Delta G^\circ = -RT\ln K = \qty{-55}{kJ/mol}$. Les six liaisons Ni--N sont de
même nature des deux côtés~; le nombre de molécules libres passe de quatre
à sept~: le gain d'entropie de translation gouverne la réaction.
\end{solution}

\begin{solution}{exo:b3:supramolecular:7}
$\Delta G^\circ = -RT\ln(\num{2.0e5}) = \qty{-30.2}{kJ/mol}$~; $T\Delta S^\circ = \Delta H^\circ - \Delta G^\circ = -40 + 30.2 = \qty{-9.8}{kJ/mol}$~;
$\Delta S^\circ = \qty{-33}{J.K^{-1}.mol^{-1}}$~: une association gouvernée par l'enthalpie, en partie payée en entropie.
\end{solution}

\begin{solution}{exo:b3:supramolecular:8}
Le cuivre(I) fixe deux motifs phénanthroline bidentés et, étant
tétraédrique, les maintient à angle droit, l'un passant à travers l'espace
où se refermera le cycle de l'autre. La fermeture des deux chaînes en
cycles les entrelace alors. Un large excès de cyanure, qui fixe bien plus
fortement le cuivre(I), retire le métal et libère le \omterm{def:b3:supramolecular:interlocked}{caténane}.
\end{solution}

\begin{solution}{exo:b3:supramolecular:9}
$k_c = \pi\Delta\nu/\sqrt2 = \pi \times 120/1.414 = \qty{267}{s^{-1}}$. Eyring~:
$\Delta G^\ddagger = RT\ln(k_BT/hk_c) = 8.314 \times 300 \times \ln(\num{6.25e12}/267) = \qty{59.6}{kJ/mol}$.
\end{solution}

\begin{solution}{exo:b3:supramolecular:10}
Avec un excès de médicament solide, le médicament libre est fixé à $S_0$. Alors $[\mathrm{DC}] = KS_0[\mathrm C]$ et
$C_t = [\mathrm C] + [\mathrm{DC}] = [\mathrm C](1 + KS_0)$, donc $[\mathrm{DC}] = KS_0C_t/(1 + KS_0)$ et la solubilité totale vaut $S_0 + [\mathrm{DC}]$. Ici
$KS_0 = 0.050$~: $S = 0.10 + 0.050 \times 10/1.050 = \qty{0.58}{mmol/L}$, près de six fois la valeur intrinsèque.
\end{solution}

\begin{solution}{exo:b3:supramolecular:11}
Lorsque $1/K \ll H_0, G_0$, $[\mathrm{HG}] \approx \bigl(H_0 + G_0 - |H_0 - G_0|\bigr)/2 = \min(H_0, G_0)$~: la fraction liée vaut $G_0/H_0$ jusqu'à un
équivalent, puis 1. La courbe ne dépend plus de $K$~: toute valeur assez
grande donne la même cassure nette aux erreurs près, si bien que les
données montrent seulement que $K$ dépasse un certain minimum.
\end{solution}

\begin{solution}{exo:b3:supramolecular:12}
$k_{\mathrm{off}} = k_{\mathrm{on}}/K = \qty{1e9}{L.mol^{-1}.s^{-1}}/\qty{1500}{L.mol^{-1}} = \qty{6.7e5}{s^{-1}}$. L'écart de déplacement
$\qty{0.160}{ppm}$ vaut \qty{64}{Hz} à \qty{400}{MHz}~; la coalescence exigerait $k \approx \pi \times 64/\sqrt2 = \qty{140}{s^{-1}}$. L'échange est
des milliers de fois plus rapide~: un seul signal moyenné.
\end{solution}

\begin{solution}{pb:b3:supramolecular:1}
\textbf{1.} Un cycle de six motifs $\ce{-CH2CH2O-}$~; dans le complexe, $\ce{K+}$ se place au centre, lié par les
six oxygènes tournés vers l'intérieur, les groupes $\ce{CH2}$ à l'extérieur.

\textbf{2.} $\ce{K+}$~: $-RT\ln 10 \times 6.1 = \qty{-34.8}{kJ/mol}$~; $\ce{Na+}$~: \qty{-25.1}{kJ/mol}.

\textbf{3.} $10^{6.1 - 4.4} = 50$.

\textbf{4.} L'ion $\ce{K+}$, plus gros (\qty{138}{pm}), atteint les six oxygènes du cycle à la fois~; l'ion
$\ce{Na+}$, plus petit (\qty{102}{pm}), ne le peut pas, sauf si le cycle se replie, ce qui coûte de la tension.

\textbf{5.} Son extérieur est une surface hydrocarbonée et sa charge est enfouie~; $\ce{KCl}$
aurait besoin que ses ions soient solvatés, ce dont le toluène est incapable.

\textbf{6.} L'eau solvate $\ce{K+}$ bien mieux que le méthanol, et l'ion doit perdre cette
eau pour entrer dans le cycle.

\textbf{7.} $\ce{K+}(\text{aq}) + \ce{MnO4-}(\text{aq}) + \mathrm L(\text{org}) \rightleftharpoons \ce{[KL]+MnO4-}(\text{org})$.

\textbf{8.} $[\ce{K+}]_{\mathrm{aq}}$ reste égale à \qty{0.10}{mol/L}~: $y = \num{1.0e5} \times 0.10 \times (\num{1.0e-3} - y)^2 = \num{1.0e4}(\num{1.0e-3} - y)^2$.

\textbf{9.} Avec $u = \num{1.0e-3} - y$~: $\num{1.0e4}u^2 + u - \num{1.0e-3} = 0$, $u = \qty{2.70e-4}{mol/L}$, $y = \qty{7.30e-4}{mol/L}$~: 73\,\%
extraits.

\textbf{10.} $y = \num{1.0e4}(\num{1.0e-3} - y)(\num{1.0e-2} - y)$ donne $y = \qty{9.89e-4}{mol/L}$~: 99\,\%.

\textbf{11.} $y = \num{1.0e4}(\num{1.0e-3} - y)(\num{1.0e-4} - y)$ donne $y = \qty{9.0e-5}{mol/L}$~: 9\,\%, limité par la couronne,
chargée à 90\,\%.

\textbf{12.} L'ion permanganate garde sa couleur violette dans le toluène.
Il y est non solvaté, anion «~nu~», et réagit bien plus vite que dans
l'eau, où une couche d'hydratation le protège.

\textbf{13.} L'ion sodium va et vient bien plus vite que la différence des
deux fréquences de résonance~: échange rapide.

\textbf{14.} $\delta = \delta_{\mathrm{free}} + (\delta_{\mathrm{bound}} - \delta_{\mathrm{free}})[\mathrm{HG}]/H_0$, avec $[\mathrm{HG}]$ tiré de l'isotherme exacte, $H_0 = \qty{2.0e-3}{mol/L}$,
$G_0 = (\text{équiv.})\,H_0$, $\delta_{\mathrm{free}} = \qty{3.560}{ppm}$.

\textbf{15.} $K = (1.55 \pm 0.08) \times 10^3\ \unit{L.mol^{-1}}$, $\delta_{\mathrm{bound}} = \qty{3.719 \pm 0.001}{ppm}$.

\textbf{16.} L'écart quadratique moyen des résidus vaut \qty{0.0014}{ppm},
de l'ordre de la précision de lecture, sans tendance~: le modèle 1:1
convient.

\textbf{17.} Fraction liée $f$ pour $G_0 = fH_0 + f/(K(1 - f))$~: 0.28 équivalent pour $f = 0.2$ et 2.1 pour $f = 0.8$.

\textbf{18.} $KH_0 = 10^{6.1} \times \num{2.0e-3} \approx 2500$~: la courbe présenterait une cassure nette à un équivalent
et ne donnerait qu'une borne inférieure~; il faut une expérience de
compétition ou une calorimétrie à plus faible concentration.

\textbf{19.} Dans le toluène, en solution homogène de l'alcène et du permanganate.

\textbf{20.} Le manganèse réduit quitte la phase organique (sous forme de dioxyde de manganèse),
et la couronne, libérée à l'interface, reprend un autre $\ce{K+}$ et un autre $\ce{MnO4-}$. Chaque molécule de
couronne transporte le permanganate de nombreuses fois~: une quantité
catalytique suffit.

\textbf{21.} $K_{\mathrm{ex}}$ multiplie $[\ce{K+}]$~: une forte concentration en potassium favorise l'extraction.

\textbf{22.} Les sels d'ammonium quaternaire, comme le chlorure de
tétrabutylammonium ou de benzyltriéthylammonium, dont le cation est
lipophile par lui-même.

\textbf{23.} Le benzène provoque des leucémies~; le toluène fait le même
travail avec un danger bien moindre.

\textbf{24.} Avec \qty{1.0}{mmol/L} de couronne, 73\,\% du permanganate est
extrait dans le toluène.
\end{solution}
